% Options for packages loaded elsewhere
\PassOptionsToPackage{unicode}{hyperref}
\PassOptionsToPackage{hyphens}{url}
\documentclass[
]{article}
\usepackage{xcolor}
\usepackage[margin=1in]{geometry}
\usepackage{amsmath,amssymb}
\setcounter{secnumdepth}{-\maxdimen} % remove section numbering
\usepackage{iftex}
\ifPDFTeX
  \usepackage[T1]{fontenc}
  \usepackage[utf8]{inputenc}
  \usepackage{textcomp} % provide euro and other symbols
\else % if luatex or xetex
  \usepackage{unicode-math} % this also loads fontspec
  \defaultfontfeatures{Scale=MatchLowercase}
  \defaultfontfeatures[\rmfamily]{Ligatures=TeX,Scale=1}
\fi
\usepackage{lmodern}
\ifPDFTeX\else
  % xetex/luatex font selection
\fi
% Use upquote if available, for straight quotes in verbatim environments
\IfFileExists{upquote.sty}{\usepackage{upquote}}{}
\IfFileExists{microtype.sty}{% use microtype if available
  \usepackage[]{microtype}
  \UseMicrotypeSet[protrusion]{basicmath} % disable protrusion for tt fonts
}{}
\makeatletter
\@ifundefined{KOMAClassName}{% if non-KOMA class
  \IfFileExists{parskip.sty}{%
    \usepackage{parskip}
  }{% else
    \setlength{\parindent}{0pt}
    \setlength{\parskip}{6pt plus 2pt minus 1pt}}
}{% if KOMA class
  \KOMAoptions{parskip=half}}
\makeatother
\usepackage{graphicx}
\makeatletter
\newsavebox\pandoc@box
\newcommand*\pandocbounded[1]{% scales image to fit in text height/width
  \sbox\pandoc@box{#1}%
  \Gscale@div\@tempa{\textheight}{\dimexpr\ht\pandoc@box+\dp\pandoc@box\relax}%
  \Gscale@div\@tempb{\linewidth}{\wd\pandoc@box}%
  \ifdim\@tempb\p@<\@tempa\p@\let\@tempa\@tempb\fi% select the smaller of both
  \ifdim\@tempa\p@<\p@\scalebox{\@tempa}{\usebox\pandoc@box}%
  \else\usebox{\pandoc@box}%
  \fi%
}
% Set default figure placement to htbp
\def\fps@figure{htbp}
\makeatother
\setlength{\emergencystretch}{3em} % prevent overfull lines
\providecommand{\tightlist}{%
  \setlength{\itemsep}{0pt}\setlength{\parskip}{0pt}}
\usepackage{bookmark}
\IfFileExists{xurl.sty}{\usepackage{xurl}}{} % add URL line breaks if available
\urlstyle{same}
% fallback for those not using the hyperref driver hyperxmp:
\makeatletter
\@ifundefined{xmpquote}{\newcommand{\xmpquote}[1]{#1}}{}
\makeatother
\hypersetup{
  pdftitle={Leverage Score and SVD},
  hidelinks,
  pdfcreator={LaTeX via pandoc}}

\title{Leverage Score and SVD}
\author{}
\date{\vspace{-2.5em}}

\begin{document}
\maketitle

\subsection{Motivation}\label{motivation}

Working with large single-cell datasets can be computationally
expensive. To reduce computational complexity, we can select a
representative subset of cells using \textbf{leverage score sampling}.

Consider the matrix:

\[
X =
\begin{pmatrix}
10 & 0 \\
10 & 0 \\
10 & 0 \\
0 & 10
\end{pmatrix}
\]

Each row represents a cell, and each column represents a gene.

The first three cells are identical, while the fourth represents a
different expression pattern.

If we want to retain only two cells, it would be reasonable to select
one of the first three cells and the fourth cell.

\textbf{Main idea:} Use Singular Value Decomposition (SVD) to identify
the contribution of each cell to the linear structure of the dataset.

\subsection{1. Singular Value
Decomposition}\label{singular-value-decomposition}

Any real matrix can be decomposed as:

\[
\boxed{X = U\Sigma V^T}
\]

where:

\begin{itemize}
\tightlist
\item
  \(U\) contains the left singular vectors.
\item
  \(\Sigma\) is a diagonal matrix containing singular values.
\item
  \(V\) contains the right singular vectors.
\end{itemize}

\subsubsection{Step 1: Computing the singular
values}\label{step-1-computing-the-singular-values}

Calculate:

\[
X^TX
\]

Find its eigenvalues:

\[
X^TXv_i=\lambda_i v_i
\]

The singular values are:

\[
\boxed{\sigma_i=\sqrt{\lambda_i}}
\]

Arrange these values in decreasing order along the diagonal of
\(\Sigma\).

\subsubsection{Step 2: Computing V}\label{step-2-computing-v}

The columns of \(V\) are the normalized eigenvectors of \(X^TX\),
ordered according to the singular values.

\[
V=(v_1,v_2,\ldots,v_r)
\]

\subsubsection{Step 3: Computing U}\label{step-3-computing-u}

Using:

\[
XV=U\Sigma
\]

we obtain:

\[
\boxed{U=XV\Sigma^{-1}}
\]

for the nonzero singular values.

\subsection{2. Computing leverage
scores}\label{computing-leverage-scores}

For the rank-\(k\) approximation, take the first \(k\) columns of \(U\).

The leverage score of cell \(i\) is:

\[
\boxed{\ell_i=\sum_{j=1}^{k}U_{ij}^{2}}
\]

This is the squared Euclidean norm of row \(i\) of \(U_k\).

A high leverage score indicates that the cell contributes strongly to
the selected singular subspace.

\subsection{3. Sampling cells}\label{sampling-cells}

We can use leverage scores to define sampling probabilities:

\[
\boxed{p_i=\frac{\ell_i}{\sum_j \ell_j}=\frac{\ell_i}{k}}
\]

Cells with higher leverage scores are more likely to be selected.

This helps preserve important directions of variation and can improve
the representation of rare cell populations.

However, rare cells are not guaranteed to have high leverage scores,
especially if their characteristic directions are excluded from the
rank-\(k\) approximation.

\end{document}
